\section{Received Prayers, Collects, Antiphons, and Hymns}\label{app:prayers}

Every prayer below reproduces an identified received text without modern
alteration or conflation. Full publication details for every Latin and English
source stand in Appendix~\ref{app:references}.

\subsection{Collects of the Season}

\begin{massbox}{First Sunday of Advent}
\textbf{Latin} (1962 Roman Missal). Excita, quaesumus, Domine, potentiam tuam,
et veni: ut ab imminentibus peccatorum nostrorum periculis, te mereamur
protegente eripi, te liberante salvari: Qui vivis et regnas cum Deo Patre in
unitate Spiritus Sancti, Deus: per omnia saecula saeculorum. \textbf{R.} Amen.

\textbf{English} (Lasance, 1945). Bestir, O Lord, Thy might, we pray Thee and
come; That, defended by Thee, we may deserve rescue from approaching dangers
brought on by our sins, and being set free by Thee, obtain our salvation. Who
livest.
\end{massbox}

\begin{massbox}{Second Sunday of Advent}
\textbf{Latin} (1962 Roman Missal). Excita, Domine, corda nostra ad praeparandas
Unigeniti tui vias: ut, per ejus adventum, purificatis tibi mentibus servire
mereamur: Qui tecum vivit et regnat in unitate.

\textbf{English} (Lefebvre, 1925). Stir up our hearts, O Lord, to prepare the
ways of Thine only-begotten Son; that through His coming we may attain to serve
Thee with purified minds. Who liveth and reigneth with Thee.
\end{massbox}

\begin{massbox}{Immaculate Conception of the Blessed Virgin Mary (8 December)}
\textbf{Latin} (1962 Roman Missal). Deus, qui per immaculatam Virginis
Conceptionem dignum Filio tuo habitaculum praeparasti: quaesumus; ut, qui ex
morte eiusdem Filii tui praevisa eam ab omni labe praeservasti, nos quoque
mundos eius intercessione ad te pervenire concedas. Per eundem Dominum.

\textbf{English} (Lasance, 1945). O God, Who, by the Immaculate Conception of
the Virgin, didst prepare a worthy dwelling for Thy Son, we beseech Thee that
Thou, Who, by the death, foreseen by Thee, of the same Thy Son, didst preserve
her from all stain, wilt grant us also, by her intercession, to come to Thee
pure in heart. Through the same.
\end{massbox}

\begin{massbox}{Third Sunday of Advent (\textit{Gaudete})}
\textbf{Latin} (1962 Roman Missal). Aurem tuam, quaesumus, Domine, precibus
nostris accommoda: et mentis nostrae tenebras gratia tuae visitationis illustra:
Qui vivis et regnas cum Deo Patre in unitate.

\textbf{English} (Lasance, 1945). Incline Thine ear to our prayers, O Lord, we
beseech Thee; and make bright the darkness of our minds by the grace of Thy
visitation. Who livest.
\end{massbox}

\begin{massbox}{Ember Wednesday of Advent (First and Second Collects)}
\textbf{First Collect --- Latin} (1962 Roman Missal). Praesta, quaesumus,
omnipotens Deus: ut redemptionis nostrae ventura solemnitas, et praesentis nobis
vitae subsidia conferat, et aeternae beatitudinis praemia largiatur. Per
Dominum.

\textbf{First Collect --- English} (Lasance, 1945). Grant us, we beseech Thee,
O Lord, that the coming solemnity of our redemption may both bring us help for
our present life and bestow on us the rewards of everlasting blessedness.
Through our Lord.

\textbf{Second Collect --- Latin} (1962 Roman Missal). Festina, quaesumus,
Domine, ne tardaveris, et auxilium nobis supernae virtutis impende: ut adventus
tui consolationibus subleventur, qui in tua pietate confidunt: Qui vivis et
regnas cum Deo Patre in unitate.

\textbf{Second Collect --- English} (Lefebvre, 1925). Hasten, we beseech Thee,
O Lord, do not tarry, and grant us the help of Thy heavenly power : that those
who trust in Thy goodness may be helped by the consolations of Thy coming : Who
livest.
\end{massbox}

\begin{massbox}{Ember Friday of Advent}
\textbf{Latin} (1962 Roman Missal). Excita, quaesumus, Domine, potentiam tuam,
et veni: ut hi, qui in tua pietate confidunt, ab omni citius adversitate
liberentur: Qui vivis et regnas cum Deo Patre in unitate.

\textbf{English} (Lasance, 1945). Bestir, O Lord, Thy might, we pray Thee, and
come; that these who trust in Thy kindness and love, may the more speedily be
delivered from all adversity. Who livest.
\end{massbox}

\begin{massbox}{Ember Saturday of Advent (Opening Collect)}
\textbf{Latin} (1962 Roman Missal). Deus, qui conspicis, quia ex nostra
pravitate affligimur: concede propitius; ut ex tua visitatione consolemur: Qui
vivis et regnas cum Deo Patre in unitate.

\textbf{English} (Lasance, 1945). O God, Who dost behold how we are afflicted by
reason of our wickedness, mercifully grant that we may be consoled by reason of
Thy visitation. Who livest.
\end{massbox}

\begin{massbox}{Fourth Sunday of Advent}
\textbf{Latin} (1962 Roman Missal). Excita, quaesumus, Domine, potentiam tuam,
et veni: et magna nobis virtute succurre; ut per auxilium gratiae tuae, quod
nostra peccata praepediunt, indulgentia tuae propitiationis acceleret: Qui vivis
et regnas cum Deo Patre in unitate.

\textbf{English} (Lasance, 1945). Bestir, O Lord, Thy might, we beseech thee,
and come; and with great power come to our aid, that, by the help of Thy grace,
that which is hindered by our sins may be hastened by Thy merciful forgiveness.
Who livest.
\end{massbox}

\begin{massbox}{Vigil of the Nativity of the Lord (24 December)}
\textbf{Latin} (1962 Roman Missal). Deus, qui nos redemptionis nostrae annua
exspectatione laetificas: praesta; ut Unigenitum tuum, quem Redemptorem laeti
suscipimus, venientem quoque judicem securi videamus, Dominum nostrum Jesum
Christum, Filium tuum: Qui tecum vivit et regnat in unitate.

\textbf{English} (Lasance, 1945). O God, Who dost gladden us year by year with
the expectation of our redemption, grant that we, who now with joy receive Thine
only begotten Son as our Redeemer, may behold Him also without fear, when He
cometh as our judge, our Lord Jesus Christ. Who with thee.
\end{massbox}

\subsection{The Seven Great O Antiphons (17--23 December)}

Latin from the Ratisbon winter Breviary (Pustet, 1888, pp.~156--157); English
from the Marquess of Bute's \textit{Roman Breviary} (Winter, 1908, p.~244).

\begin{massbox}{17 December: \textit{O Sapientia}}
\textbf{Latin.} O Sapientia, quae ex ore Altissimi prodiisti, attingens a fine
usque ad finem, fortiter suaviterque disponens omnia: veni ad docendum nos viam
prudentiae.

\textbf{English.} O Wisdom, That comest out of the mouth of the Most High, That
reachest from one end to another, and dost mightily and sweetly order all
things: come, to teach us the way of prudence!
\end{massbox}

\begin{massbox}{18 December: \textit{O Adonai}}
\textbf{Latin.} O Adonai, et Dux domus Isra\"el, qui Moysi in igne flammae rubi
apparuisti, et ei in Sina legem dedisti: veni ad redimendum nos in brachio
extento.

\textbf{English.} O Adonai, and Ruler of the house of Isra\"el, Who didst appear
unto Moses in the burning bush, and gavest him the law in Sinai: come, to redeem
us with an outstretched arm!
\end{massbox}

\begin{massbox}{19 December: \textit{O Radix Jesse}}
\textbf{Latin.} O radix Jesse, qui stas in signum populorum, super quem
continebunt reges os suum, quem Gentes deprecabuntur: veni ad liberandum nos,
jam noli tardare.

\textbf{English.} O Root of Jesse, Which standest for an ensign of the people,
at Whom the kings shall shut their mouths, unto Whom the Gentiles shall seek:
come, to deliver us, make no tarrying!
\end{massbox}

\begin{massbox}{20 December: \textit{O Clavis David}}
\textbf{Latin.} O clavis David, et sceptrum domus Isra\"el; qui aperis, et nemo
claudit; claudis, et nemo aperit: veni, et educ vinctum de domo carceris,
sedentem in tenebris, et umbra mortis.

\textbf{English.} O Key of David, and Sceptre of the house of Isra\"el; That
openest, and no man shutteth; and shuttest and no man openeth: come, to bring
out the prisoners from the prison, and them that sit in darkness, and in the
shadow of death!
\end{massbox}

\begin{massbox}{21 December: \textit{O Oriens}}
\textbf{Latin.} O Oriens, splendor lucis aeternae, et sol justitiae: veni, et
illumina sedentes in tenebris, et umbra mortis.

\textbf{English.} O Day-Spring, Brightness of the everlasting Light, Sun of
Righteousness, come, to give light to them that sit in darkness, and in the
shadow of death!
\end{massbox}

\begin{massbox}{22 December: \textit{O Rex Gentium}}
\textbf{Latin.} O Rex Gentium, et desideratus earum, lapisque angularis, qui
facis utraque unum: veni, et salva hominem, quem de limo formasti.

\textbf{English.} O King of the Gentiles, yea, and Desire thereof, O
Corner-stone that makest of twain one: come, to save man, whom Thou hast made of
the dust of the earth!
\end{massbox}

\begin{massbox}{23 December: \textit{O Emmanuel}}
\textbf{Latin.} O Emmanuel, Rex et legifer noster, exspectatio Gentium, et
Salvator earum: veni ad salvandum nos Domine Deus noster.

\textbf{English.} O Emmanuel, our King and our Law-giver, Longing of the
Gentiles, yea, and Salvation thereof: come to save us, O Lord our God!
\end{massbox}

\subsection{Seasonal Hymns, Marian Antiphon, and Devotions}

\begin{massbox}{Vespers Hymn of Advent: \textit{Creator alme siderum} (Latin)}
\textbf{Latin} (Pustet, 1888, p.~132).\\
Creator alme siderum,\\
Aeterna lux credentium,\\
Jesu Redemptor omnium,\\
Intende votis supplicum.\\[0.4em]
Qui daemonis ne fraudibus\\
Periret orbis, impetu\\
Amoris actus, languidi\\
Mundi medela factus es.\\[0.4em]
Commune qui mundi nefas\\
Ut expiares; ad crucem\\
E Virginis sacrario\\
Intacta prodis victima.\\[0.4em]
Cujus potestas gloriae,\\
Nomenque cum primum sonat;\\
Et coelites et inferi\\
Tremente curvantur genu.\\[0.4em]
Te deprecamur ultimae\\
Magnum diei Judicem,\\
Armis supernae gratiae\\
Defende nos ab hostibus.\\[0.4em]
Virtus, honor, laus, gloria\\
Deo Patri cum Filio,\\
Sancto simul Paraclito,\\
In saeculorum saecula. Amen.
\end{massbox}

\begin{massbox}{Vespers Hymn of Advent: \textit{Creator alme siderum} (English)}
\textbf{English} (St.~John Henry Newman, in Bute, 1908, p.~213).\\
Creator of the starry pole,\\
Saviour of all who live,\\
And light of every faithful soul,\\
JESU, these prayers receive.\\[0.4em]
Who sooner than our foe malign\\
Should triumph, from above\\
Didst come, to be the medicine\\
Of a sick world, in love;\\[0.4em]
And the deep wounds to cleanse and cure\\
Of a whole race, didst go,\\
Pure Victim, from a Virgin pure,\\
The bitter Cross unto.\\[0.4em]
Who hast a Name, and hast a Power,\\
The height and depth to sway,\\
And Angels bow, and devils cower,\\
In transport or dismay;\\[0.4em]
Thou too shalt be our Judge at length;\\
Lord, in Thy grace bestow\\
Thy weapons of celestial strength,\\
And snatch us from the foe.\\[0.4em]
Honour and glory, power and praise,\\
To Father, and to Son,\\
And Holy Ghost, be paid always,\\
The Eternal Three in One. Amen.
\end{massbox}

\begin{massbox}{Marian Antiphon of Advent: \textit{Alma Redemptoris Mater}}
\textbf{Latin} (Pustet, 1888, pp.~130--131; from First Vespers of Advent~I
through Compline of 2 February). Alma Redemptoris Mater, quae pervia coeli Porta
manes, et stella maris, succurre cadenti, Surgere qui curat, populo: tu quae
genuisti, Natura mirante, tuum sanctum Genitorem, Virgo prius ac posterius,
Gabrielis ab ore Sumens illud Ave, peccatorum miserere.

\textbf{In Advent:} \textbf{V.} Angelus Domini nuntiavit Mariae. \textbf{R.} Et
concepit de Spiritu sancto. \textit{Oremus.} Gratiam tuam, quaesumus Domine,
mentibus nostris infunde: ut qui, Angelo nuntiante, Christi Filii tui
Incarnationem cognovimus, per passionem ejus et crucem ad resurrectionis gloriam
perducamur. Per eumdem Christum Dominum nostrum. \textbf{R.} Amen.

\textbf{English} (Bute, 1908, pp.~210--211). Maiden! Mother of Him Who redeemed
us, thou that abidest Heaven's open gate, and the Star of the Sea, come, succour
the fallen! Fallen indeed we are, but fain would rise by thy succour. Thou that
beyond nature's course hast borne in time the Eternal; Thou that a Virgin
before, and after that childbirth remainest, From the Archangel's lips the
quickening message receiving, Mother of JESUS and us, turn thine eyes of mercy on
sinners.

\textbf{In Advent:} \textbf{Verse.} The Angel of the Lord announced unto Mary,
\textbf{Answer.} And she conceived of the Holy Ghost. \textit{Let us pray.} We
beseech Thee, O Lord, pour Thy grace into our hearts; that, as we have known the
Incarnation of Thy Son JESUS Christ by the message of an Angel, so by His
Passion and Cross we may be brought unto the glory of His Resurrection. Through
the same Christ our Lord. \textbf{Answer.} Amen.

\textit{From First Vespers of Christmas onward, the versicle changes to
``Post partum Virgo inviolata permansisti / Dei Genitrix intercede pro nobis''
(``After thy delivery thou still remainest a Virgin undefiled / Mother of God,
pray for us'') and the Collect changes to ``Deus, qui salutis aeternae.''}
\end{massbox}

\begin{massbox}{Advent Prose: \textit{Rorate caeli} (Latin)}
\textbf{Latin} (Solesmes, \textit{Manual of Gregorian Chant}, 1903, pp.~313--315;
a French Advent chant and Benediction prose, not a Roman Breviary text).
Rorate caeli desuper, et nubes pluant justum.\\[0.3em]
1. Ne irascaris Domine, ne ultra memineris iniquitatis: ecce civitas Sancti
facta est deserta: Sion deserta facta est: Jerusalem desolata est: domus
sanctificationis tuae et gloriae tuae, ubi laudaverunt te patres nostri.
\textbf{R.} Rorate.\\[0.3em]
2. Peccavimus, et facti sumus tamquam immundus nos, et cecidimus quasi folium
universi: et iniquitates nostrae quasi ventus abstulerunt nos: abscondisti
faciem tuam a nobis, et allisisti nos in manu iniquitatis nostrae. \textbf{R.}
Rorate.\\[0.3em]
3. Vide Domine afflictionem populi tui, et mitte quem missurus es: emitte Agnum
dominatorem terrae, de petra deserti ad montem filiae Sion: ut auferat ipse
jugum captivitatis nostrae. \textbf{R.} Rorate.\\[0.3em]
4. Consolamini, consolamini popule meus: cito veniet salus tua: quare moerore
consumeris, quia innovavit te dolor? Salvabo te, noli timere, ego enim sum
Dominus Deus tuus, Sanctus Israel, redemptor tuus. \textbf{R.} Rorate.
\end{massbox}

\begin{massbox}{Advent Prose: \textit{Rorate caeli} (English)}
\textbf{English} (Dom Prosper Gu\'eranger / Dom Laurence Shepherd, 1870,
pp.~155--156, translating a Latin text with minor variants in verses 1, 3, and
4). Drop down dew, ye heavens, from above, and let the clouds rain the Just
One.\\[0.3em]
1. Be not angry, O Lord, and remember no longer our iniquity: behold the city of
thy Sanctuary is become a desert, Sion is made a desert. Jerusalem is desolate,
the house of our holiness and of thy glory, where our fathers praised thee.
\textit{[Refrain]}\\[0.3em]
2. We have sinned, and we are become as one unclean, and we have all fallen as a
leaf; and our iniquities, like the wind, have taken us away: thou hast hid thy
face from us, and hast crushed us by the hand of our iniquity.
\textit{[Refrain]}\\[0.3em]
3. See, O Lord, the affliction of thy people, and send Him whom thou hast
promised to send. Send forth the Lamb, the ruler of the earth, from the rock of
the desert to the mount of the daughter of Sion, that himself may take off the
yoke of our captivity. \textit{[Refrain]}\\[0.3em]
4. Be comforted, be comforted, my people; thy salvation shall speedily come: why
wilt thou waste away in sadness? why hath sorrow seized thee? I will save thee;
fear not: for I am the Lord thy God, the Holy One of Israel, thy Redeemer. Drop
down dew, ye heavens, from above, and let the clouds rain the Just one.
\end{massbox}

\begin{massbox}{The Angelus}
\textbf{Latin} (\textit{Raccolta}, 1910, pp.~151--152). \textbf{V.} Angelus
DOMINI nuntiavit Mariae. \textbf{R.} Et concepit de SPIRITU SANCTO. \textit{Ave
Maria.}\\
\textbf{V.} Ecce ancilla DOMINI. \textbf{R.} Fiat mihi secundum verbum tuum.
\textit{Ave Maria.}\\
\textbf{V.} Et Verbum caro factum est. \textbf{R.} Et habitavit in nobis.
\textit{Ave Maria.}\\
\textbf{V.} Ora pro nobis, sancta DEI Genitrix. \textbf{R.} Ut digni efficiamur
promissionibus CHRISTI.\\
\textit{Oremus.} Gratiam tuam, quaesumus, DOMINE, mentibus nostris infunde: ut
qui, angelo nuntiante, CHRISTI FILII tui Incarnationem cognovimus, per Passionem
ejus et Crucem ad resurrectionis gloriam perducamur. Per eundem CHRISTUM DOMINUM
nostrum. \textbf{R.} Amen.

\textbf{English} (Ambrose St.~John, \textit{Raccolta}, 1910, pp.~151--152).
\textbf{V.} The angel of the LORD declared unto Mary. \textbf{R.} And she
conceived of the HOLY GHOST. \textit{Ave Maria.}\\
\textbf{V.} Behold the handmaid of the LORD. \textbf{R.} Be it done unto me
according to thy word. \textit{Ave Maria.}\\
\textbf{V.} And the Word was made Flesh. \textbf{R.} And dwelt among us.
\textit{Ave Maria.}\\
\textbf{V.} Pray for us, O Holy Mother of GOD. \textbf{R.} That we may be made
worthy of the promises of CHRIST.\\
\textit{Let us pray.} Pour forth, we beseech Thee, O LORD, thy grace into our
hearts; that we, to whom the Incarnation of CHRIST thy SON was made known by the
message of an angel, may by his Passion and Cross be brought to the glory of his
resurrection. Through the same CHRIST our LORD. \textbf{R.} Amen.
\end{massbox}

\subsection{Gospel Canticles of the Divine Office}

Latin from the Clementine Vulgate collated against liturgical editions
(correcting the electronic \texttt{latVUC} error at Luke~1:70 to
\textit{locutus}; see Appendix~\ref{app:about}); English from the American
Douay-Rheims Bible (1899). In the Divine Office each canticle closes with the
\textit{Gloria Patri}.

\begin{massbox}{\textit{Magnificat} (Vespers, Luke 1:46--55)}
\textbf{Latin.} Magnificat anima mea Dominum: et exsultavit spiritus meus in Deo
salutari meo. Quia respexit humilitatem ancill\ae{} su\ae{}: ecce enim ex hoc
beatam me dicent omnes generationes, quia fecit mihi magna qui potens est: et
sanctum nomen ejus, et misericordia ejus a progenie in progenies timentibus
eum. Fecit potentiam in brachio suo: dispersit superbos mente cordis sui.
Deposuit potentes de sede, et exaltavit humiles. Esurientes implevit bonis: et
divites dimisit inanes. Suscepit Isra\"el puerum suum, recordatus
misericordi\ae{} su\ae{}: sicut locutus est ad patres nostros, Abraham et semini
ejus in s\ae cula.

\textbf{English.} My soul doth magnify the Lord. And my spirit hath rejoiced in
God my Saviour. Because he hath regarded the humility of his handmaid; for
behold from henceforth all generations shall call me blessed. Because he that is
mighty, hath done great things to me; and holy is his name. And his mercy is
from generation unto generations, to them that fear him. He hath shewed might in
his arm: he hath scattered the proud in the conceit of their heart. He hath put
down the mighty from their seat, and hath exalted the humble. He hath filled the
hungry with good things; and the rich he hath sent empty away. He hath received
Israel his servant, being mindful of his mercy: As he spoke to our fathers, to
Abraham and to his seed for ever.
\end{massbox}

\begin{massbox}{\textit{Benedictus} (Lauds, Luke 1:68--79)}
\textbf{Latin.} Benedictus Dominus Deus Isra\"el, quia visitavit, et fecit
redemptionem plebis su\ae{}: et erexit cornu salutis nobis in domo David pueri
sui, sicut locutus est per os sanctorum, qui a s\ae culo sunt, prophetarum ejus:
salutem ex inimicis nostris, et de manu omnium qui oderunt nos: ad faciendam
misericordiam cum patribus nostris: et memorari testamenti sui sancti:
jusjurandum, quod juravit ad Abraham patrem nostrum, daturum se nobis ut sine
timore, de manu inimicorum nostrorum liberati, serviamus illi in sanctitate et
justitia coram ipso, omnibus diebus nostris. Et tu puer, propheta Altissimi
vocaberis: pr\ae ibis enim ante faciem Domini parare vias ejus, ad dandam
scientiam salutis plebi ejus in remissionem peccatorum eorum per viscera
misericordi\ae{} Dei nostri, in quibus visitavit nos, oriens ex alto: illuminare
his qui in tenebris et in umbra mortis sedent: ad dirigendos pedes nostros in
viam pacis.

\textbf{English.} Blessed be the Lord God of Israel; because he hath visited and
wrought the redemption of his people: And hath raised up an horn of salvation to
us, in the house of David his servant: As he spoke by the mouth of his holy
prophets, who are from the beginning: Salvation from our enemies, and from the
hand of all that hate us: To perform mercy to our fathers, and to remember his
holy testament, The oath, which he swore to Abraham our father, that he would
grant to us, That being delivered from the hand of our enemies, we may serve him
without fear, In holiness and justice before him, all our days. And thou, child,
shalt be called the prophet of the Highest: for thou shalt go before the face of
the Lord to prepare his ways: To give knowledge of salvation to his people, unto
the remission of their sins: Through the bowels of the mercy of our God, in
which the Orient from on high hath visited us: To enlighten them that sit in
darkness, and in the shadow of death: to direct our feet into the way of peace.
\end{massbox}

\begin{massbox}{\textit{Nunc Dimittis} (Compline, Luke 2:29--32)}
\textbf{Latin.} Nunc dimittis servum tuum Domine, secundum verbum tuum in pace:
quia viderunt oculi mei salutare tuum, quod parasti ante faciem omnium
populorum: lumen ad revelationem gentium, et gloriam plebis tu\ae{} Isra\"el.

\textbf{English.} Now thou dost dismiss thy servant, O Lord, according to thy
word in peace; Because my eyes have seen thy salvation, Which thou hast prepared
before the face of all peoples: A light to the revelation of the Gentiles, and
the glory of thy people Israel.
\end{massbox}
